% Fig. 3 -- case-study timelines
% Included by paper/icc_main.tex via \input.
% Node sublabels name the satellite-specific object each stage owns.
\begin{figure}[!t]
\centering
\resizebox{0.94\columnwidth}{!}{%
\begin{tikzpicture}[font=\scriptsize,>={Latex[length=2.0mm]},
  nd/.style={draw=black!70,line width=0.6pt,rounded corners=1.5pt,align=center,
             minimum height=9mm,minimum width=20mm,inner sep=2pt,fill=white},
  fl/.style={->,line width=0.8pt,black!70},
  bad/.style={->,line width=0.8pt,red!65!black,dashed},
  stop/.style={draw=red!65!black,line width=1.1pt,rounded corners=2pt,fill=red!6,
               align=center,inner sep=2.5pt,font=\scriptsize}]
% ---------------- panel (a)
\node[anchor=west,font=\scriptsize\bfseries] at (-0.3,2.05) {(a) retrospective-label pipeline};
\node[nd] (a1) at (0.9,1.05) {held element\\[-1pt]{\tiny published $\le t_d$}};
\node[nd] (a2) at (3.6,1.05) {visible-pass\\[-1pt]{\tiny above mask}};
\node[nd] (a3) at (6.3,1.05) {frozen\\[-1pt]{\tiny registry hashed}};
\node[nd] (a4) at (9.0,1.05) {later element\\[-1pt]{\tiny catalogue publishes}};
\node[nd] (a5) at (11.7,1.05) {label closes\\[-1pt]{\tiny at $t_c$}};
\foreach \x/\y in {a1/a2,a2/a3,a3/a4,a4/a5} \draw[fl] (\x) -- (\y);
\draw[bad] (a4.north) to[out=140,in=40]
  node[pos=0.5,above,font=\scriptsize,fill=white,inner sep=1pt]
  {element-epoch gap as a feature} (a3.north);
\draw[bad] (a5.south) to[out=-150,in=-30]
  node[pos=0.5,below,font=\scriptsize,fill=white,inner sep=1pt]
  {row membership set by later publication} (a3.south);
\node[stop] at (6.3,-1.0) {\textsc{halt} \rulename{L1.1} \rulename{L1.4}};
\draw[red!65!black,line width=0.9pt] (6.3,0.6) -- (6.3,-0.62);
% ---------------- panel (b)
\node[anchor=west,font=\scriptsize\bfseries] at (-0.3,-1.95) {(b) learning and gating pipeline};
\node[nd] (b1) at (0.9,-3.0) {physics\\[-1pt]{\tiny baseline}};
\node[nd] (b2) at (3.6,-3.0) {learned\\[-1pt]{\tiny residual branch}};
\node[nd] (b3) at (6.3,-3.0) {validate,\\[-1pt]{\tiny select}};
\node[nd] (b4) at (9.0,-3.0) {\textbf{freeze}\\[-1pt]{\tiny model, gate, tracker}};
\node[nd] (b5) at (11.7,-3.0) {deployment\\[-1pt]{\tiny per pass}};
\foreach \x/\y in {b1/b2,b2/b3,b3/b4,b4/b5} \draw[fl] (\x) -- (\y);
\draw[bad] (b5.north) to[out=140,in=40]
  node[pos=0.5,above,font=\scriptsize,fill=white,inner sep=1pt]
  {tracker state advances past the freeze} (b3.north);
\node[font=\scriptsize,text=red!65!black,anchor=west] at (0.15,-4.15)
  {zero-effect control admits $\Rightarrow$ target not null};
\node[font=\scriptsize,text=red!65!black,anchor=west] at (0.15,-4.62)
  {pooled per sample, not per pass $\Rightarrow$ unit finer than exchangeable};
\node[stop] at (9.0,-5.35) {\textsc{halt} \rulename{L3.1} \rulename{L3.3}
  \rulename{L4.2} \rulename{L4.7}};
\draw[red!65!black,line width=0.9pt] (9.0,-3.45) -- (9.0,-5.0);
\end{tikzpicture}}
\caption{Where the contract halts each case study, with the satellite object each stage owns
named beneath it: element publication versus epoch, visible-pass scheduling, the frozen
registry fixing row membership, frozen model/gate/tracker state, and the pass as statistical
unit. Solid arrows are the intended flow, dashed the defects found. Every dashed path is
chronologically consistent, which is why ordering does not exclude it. No model-performance
quantity is reported from either pipeline.}
\label{fig:cases}
\end{figure}
